% !TeX program = xelatex
\documentclass[11pt,a4paper,fontset=fandol]{ctexart}
\usepackage[sexy,noasy,noauthor,nofancy]{evan}
\usepackage{geometry,booktabs,tikz,fancyhdr}
\geometry{margin=2cm,headheight=15pt}
\renewcommand{\proofname}{证明}
\numberwithin{equation}{section}
\ctexset{section={format=\Large\bfseries\raggedright},subsection={format=\large\bfseries\raggedright}}
\setlength{\parskip}{.25em}
\linespread{1.08}
\pagestyle{fancy}\fancyhf{}
\fancyhead[L]{\small\nouppercase{\leftmark}}
\fancyfoot[C]{\thepage}
\hypersetup{linkcolor=blue!45!black,urlcolor=blue!45!black}
\title{数学讲义标题}
\author{LeyuDame}\date{}
\begin{document}
\maketitle
\section*{引言}
从具体的数学对象与问题引入本节。
\section{基本概念与主要结论}
\begin{definition}
给出所需对象的定义及适用范围。
\end{definition}
\begin{theorem}
列明假设，陈述结论。
\end{theorem}
\begin{proof}
写出完整且可核查的推导。
\end{proof}
\begin{example}
给出应用上述结论的具体问题。
\end{example}
\begin{proof}[解]
计算并解释结果。
\end{proof}
\section{习题}
\begin{enumerate}
\item 给出一个检验结论适用条件的问题。
\end{enumerate}
\appendix
\section{习题参考解答}
给出简洁的证明或计算。
\end{document}
